\chapter{Skill Composition and Workflows}

Skills are predefined workflows. Composing them effectively requires understanding their scope and interactions.

\section{Skill Composition Patterns}

\begin{table}[H]
\centering
\begin{tabular}{lll}
\toprule
\textbf{Pattern} & \textbf{Example} & \textbf{DCF Checkpoint} \\
\midrule
Sequential & Code → /commit → Push & Review commit message before push \\
Conditional & Test pass? → /commit : Fix & Decide whether to proceed \\
Nested & /feature-dev includes /commit & Review feature before commit phase \\
\bottomrule
\end{tabular}
\caption{Skill Composition Patterns}
\end{table}

\section{When NOT to Use Skills}

\begin{itemize}
    \item When you need fine-grained control
    \item When the skill's assumptions don't match your context
    \item When you want to learn the underlying process
\end{itemize}

\section{Creating Skills from Effective Patterns}\label{sec:pattern-to-skill}\index{Skill Creation}\index{Pattern Capture}

When you discover a prompting pattern that works particularly well, \textbf{codify it as a skill}. This transforms tacit knowledge into reusable infrastructure.

\subsection{Recognizing Skill-Worthy Patterns}

A pattern is worth capturing as a skill when:
\begin{itemize}
    \item You've used it successfully multiple times
    \item It applies to a category of tasks, not just one instance
    \item The sequence of prompts or approach is non-obvious
    \item You find yourself explaining it to others or re-discovering it
\end{itemize}

\subsection{The Pattern-to-Skill Process}

\begin{enumerate}
    \item \textbf{Notice repetition}: You're doing the same sequence again
    \item \textbf{Extract the essence}: What's the core approach that makes this work?
    \item \textbf{Generalize}: Remove instance-specific details, keep the transferable structure
    \item \textbf{Document as skill}: Create a \texttt{.claude/skills/skillname.md} file
    \item \textbf{Test and refine}: Use it on new instances, adjust based on results
\end{enumerate}

\subsection{Skill File Structure}

\begin{lstlisting}
# Skill Name

Brief description of what this skill does.

## Usage

```
/skillname              # Basic invocation
/skillname <args>       # With arguments
```

## Instructions

[Detailed guidance for Claude on how to execute this skill]

### When to Apply
[Criteria for when this skill is appropriate]

### Process
[Step-by-step approach]

### Output Format
[Expected structure of results]
\end{lstlisting}

\subsection{Example: The DCF Skill Itself}

The \texttt{/dcf} skill in this repository is itself an example of pattern capture. The underlying pattern---Socratic questioning applied to AI collaboration---was discovered through practice, recognized as reusable, and codified into a principle-based skill with 24 modes across five categories: evaluation \& review (\texttt{/dcf review}, \texttt{/dcf checkpoint}, \texttt{/dcf self-review}, \texttt{/dcf refine}), problem solving (\texttt{/dcf debug}, \texttt{/dcf unstick}, \texttt{/dcf simplify}, \texttt{/dcf diagnose}, \texttt{/dcf decompose}, \texttt{/dcf verify}), design \& analysis (\texttt{/dcf architect}, \texttt{/dcf tradeoffs}, \texttt{/dcf assumptions}, \texttt{/dcf premortem}, \texttt{/dcf challenge}, \texttt{/dcf decide}, \texttt{/dcf constrain}), learning \& exploration (\texttt{/dcf learn}, \texttt{/dcf onboard}, \texttt{/dcf explain}), and session management (\texttt{/dcf compact}, \texttt{/dcf context-health}, \texttt{/dcf retro}, \texttt{/dcf skill}). The skill is \emph{principle-based}---each mode describes an outcome, not a script, trusting Claude to apply Socratic questioning contextually. These modes can be chained into workflows for common scenarios---see Workflow Composition in the Glossary.

\subsection{DCF Connection}

Creating skills from patterns embodies multiple DCF principles:
\begin{itemize}
    \item \textbf{Language as infrastructure}: Skills are crystallized prompting knowledge
    \item \textbf{Scaffolding}: Skills scaffold future interactions with proven approaches
    \item \textbf{PKM integration}: Your skill library is a form of externalized expertise
    \item \textbf{Recursive refinement}: Skills improve through use and iteration
\end{itemize}

The act of creating a skill also forces \textbf{articulation}---you must understand the pattern well enough to explain it, which deepens your own mastery.

